\section{不确定性原理}
关于不确定性原理，受限于时间，我们仅简要介绍一下结论。微观世界的工作方式与宏观世界是不同的，例如，宏观世界中我们可以很容易的确定某个物体的位置和动量，微观世界中
\begin{itemize}
    \item 如果位置确定了，那么动量就非常不确定。
    \item 如果动量确定了，那么位置就非常不确定。
\end{itemize}
简而言之，我们无法同时确定微观粒子的位置和动量，这种想法可以表达为数学式。
\begin{BoxPrinciple}[不确定性原理]
    微观粒子的位置和动量无法同时被确定
    \begin{Equation}
        \delt{x}\cdot\delt{p}\geq\frac{\hbar}{2}
    \end{Equation}
    其中$\hbar$称为\uwave{约化普朗克常数}，其为
    \begin{Equation}
        \hbar=\frac{h}{2\pi}
    \end{Equation}
\end{BoxPrinciple}
值得注意的是，位置和动量存在不确定性关系，实际上，能量和时间之间也有不确定性原理
\begin{Equation}
    \delt{E}\cdot\delt{t}\geq\frac{\hbar}{2}
\end{Equation}
如果试图用电子的单缝衍射讨论不确定性原理，得到的是比$\delt{x}\cdot\delt{p}\geq\hbar/2$弱的$\delt{x}\cdot\delt{p}\geq h$。